\begin{claim}{D.2.49}{Proved}{The universal composite has zero integral loop Hurewicz}
Put $X_0=\Sigma B$, $Y=\Sigma X_0$, and $L=Y\langle7\rangle$.
Retain $e=u_4E\epsilon'$ and $h=h_Ml_4$ from D.2.46.
There is an actual class $t\in\pi_{10}X_0$ such that
\[
                    e=(Et)\nu_{11}\quad\text{in }\pi_{14}Y.
                    \tag{D.2.49a}
\]
The map $Et:S^{11}\to Y$ lifts uniquely through $L$ up to
homotopy, and this factorization gives the unique lift of $e$.
For the unique lift $\widetilde h\in\pi_{14}L$ one has
\[
 h_{\ZZ}\bigl(\operatorname{adj}(\widetilde h)\bigr)=0
                  \quad\text{in }H_{13}(\Omega L;\ZZ).
 \tag{D.2.49b}
\]
Consequently, if $a_6$ represents the specified integral bottom
generator in a Pontryagin chain model for $\Omega L$ and
$db_7=4a_6$, then
\[
 z_L=[a_6b_7-b_7a_6]\in H_{13}(\Omega L;\ZZ)
                         \quad\text{satisfies}\quad 2z_L=0.
 \tag{D.2.49c}
\]
The doubled class is independent of these chain choices.
These statements do not assert $h=0$ or evaluate the order of $q$
in D.2.46. The integral cohomology consequence is derived in D.2.50.
\end{claim}
\begin{proof}
\textbf{An actual injective low suspension.}
Write $u_3:S^3\to X_0$ for the bottom map, so $Eu_3=u_4$.
We first verify
\[
 \pi_6X_0=\ZZ/2\{\beta_0\},\qquad
 E\beta_0=2\alpha\in\pi_7Y=\ZZ/4\{\alpha\}.
 \tag{D.2.49d}
\]
The first Ganea evaluation construction for
$B=\Omega K(\FF_2,3)$ gives the actual fibration
\[
 G=\Sigma(B\wedge B)\longrightarrow\Sigma B
                    \longrightarrow K(\FF_2,3).
\]
Its fiber is the join $B*B$: pulling the two cone pieces
back to the path-space model gives the homotopy pushout of
the two projections $B\times B\rightrightarrows B$.
Thus $\pi_6G\to\pi_6X_0$ is an isomorphism.

Here is the complete low attaching calculation for $G$.
Use the standard classes $a_2,b_3,c_4$ of $B$, with
$\operatorname{Sq}_*^1b=a$ and
$\operatorname{Sq}_*^2c=a$.
The mod-two groups of $G$ in degrees five, six and seven
have respective bases
\[
                aa;\qquad ab,ba;\qquad ac,ca,bb.
\]
The first Bockstein takes both $ab,ba$ to $aa$, and takes
$bb$ to $ab+ba$.  The second homology square takes each
of $ac,ca,bb$ to $aa$.
In integral cellular coordinates put $B_6=ab$ and
$C_6=ba-ab$, so $dB_6=2aa$ and $dC_6=0$.
The six-skeleton is $M(\ZZ/2,5)\vee S^6_C$ and
\[
 \pi_6(G^{(6)})=\ZZ/2\{i\eta_5\}\oplus\ZZ\{C_6\}.
 \tag{D.2.49e}
\]
The actual integral cycles $ac,ca$ have zero $C_6$ degree;
their square detector therefore attaches each by $i\eta_5$.
The $bb$ cell attaches by $\pm2C_6+i\eta_5$.
These are all components of (D.2.49e), which is already in
the stable range.  Higher cells cannot change $\pi_6$.
It follows that $\pi_6G=\ZZ/2$, and its integral Hurewicz
map to $H_6G=\ZZ/2$ is an isomorphism.

In the first James pair $(JX_0,X_0)$, relative Hurewicz gives
\[
 \pi_6(JX_0,X_0)=H_6(JX_0,X_0;\ZZ)
       =H_3(X_0;\ZZ)\otimes H_3(X_0;\ZZ)=\ZZ/2.
 \tag{D.2.49f}
\]
This is the first relative word of length two in the reduced
product filtration \cite[Sections 3 and 7]{I_James1955}.
The class $\alpha=u_4\nu_4$ of D.2.18 maps to its nonzero
element: naturality from the bottom $S^3$ identifies the
relative map with reduction of the Hopf invariant one of
$\nu_4$.  Exactness, together with
$\pi_7Y=\ZZ/4\{\alpha\}$, makes the image of suspension
exactly $\langle2\alpha\rangle$.
The preceding computation of its source proves (D.2.49d).
In particular this suspension is injective.

\textbf{The full Toda indeterminacy supplies a sphere factor.}
D.2.44 proved the actual relation $u_4E\nu'=0$.
Equation (D.2.49d) now implies
\[
                         u_3\nu'=0\quad\text{in }\pi_6X_0.
                         \tag{D.2.49g}
\]
The exact sphere input is
\[
 \epsilon'\in\{\nu',2\nu_6,\nu_9\}_3
                  \subseteq\{\nu',2\nu_6,\nu_9\}.
 \tag{D.2.49h}
\]
The indexed definition, its convention, and the same
zero-first-entry argument are explicit in the opened author
source \cite[Section 2.5, p.~4, and Proposition 2, p.~6]
{FGHStricklandToda}; its attribution is Toda, Section VI(iv).
Only membership in the ordinary bracket is needed here.
Postcomposition and (D.2.49g) give
\[
 u_3\epsilon'\in\{0,2\nu_6,\nu_9\}
          =\pi_{10}X_0\circ\nu_{10}.
 \tag{D.2.49i}
\]
The equality includes every choice of nullhomotopy.
With zero first map, the bracket contains zero, its first
indeterminacy summand is zero, and its entire other summand
is $[S^{10},X_0]\circ\Sigma\nu_9$.
Thus there is an actual $t$ with $u_3\epsilon'=t\nu_{10}$;
suspending proves (D.2.49a).  No numerical value of this
remaining composition is inferred.

\textbf{Hurewicz vanishing in the connected cover.}
For $j>6$, the map $\pi_jL\to\pi_jY$ is an isomorphism.
In particular $Et$ and its homotopies lift uniquely through
$L$, and the lifted factorization in (D.2.49a) is the unique
lift of $e$.  Its once-adjoint factors through $\Omega S^{11}$.
The once-adjoint of the lift of $\alpha\sigma'$ similarly
factors through $\Omega S^7$.
The integral James tensor calculation gives
\[
 H_{13}(\Omega S^{11};\ZZ)=0,
                 \qquad H_{13}(\Omega S^7;\ZZ)=0,
 \tag{D.2.49j}
\]
since the respective tensor generators have degrees ten and
six.  Hence both summands of the lift of
$q=\alpha\sigma'+e$ have zero integral loop Hurewicz.
D.2.46 places $h$ in $\langle q\rangle$; the isomorphism
on $\pi_{14}$ and the Hurewicz homomorphism therefore prove
(D.2.49b).  The factor through $S^{11}$ is essential:
factoring $e$ through $S^4$ only gives a conclusion in
$\Omega Y$, because its bottom map does not lift to $L$.

\medskip\noindent
\textbf{The intrinsic doubled commutator.}
Let $M_4=M(\mathbb Z/4,7)$ and use a Pontryagin differential
graded algebra for the actual loop space $\Omega L$.  Choose a
cycle $a$ of degree six representing the specified bottom
generator and a chain $b$ of degree seven with $db=4a$.
The low homology calculation in D.2.22, after the two homology
suspensions to $L$, gives
\[
 H_6(\Omega L;\mathbb Z)=\mathbb Z/4,
 \qquad H_7(\Omega L;\mathbb Z)=\mathbb Z/2.
 \tag{D.2.49e1}
\]
These are homology groups; the integral cohomology group
$H^8(L;\mathbb Z)$ in D.2.47--48 is instead $\mathbb Z/4$.
Put $z_L=[ab-ba]\in H_{13}(\Omega L;\mathbb Z)$.  We prove
the exact comparison
\[
 2z_L=h_{\mathbb Z}\bigl(\operatorname{adj}(h_Ml_4)\bigr).
 \tag{D.2.49e2}
\]

First this doubled class is intrinsic.  The identity
$d(b^2)=4(ab-ba)$ gives $4z_L=0$.  If $a'=a+dv$, with
$|v|=7$, take $b'=b+4v$.  The graded Leibniz rule gives the
literal identity
\[
 a'b'-b'a'-(ab-ba)=d(vb+bv+4v^2).
 \tag{D.2.49e3}
\]
Any further change of $b'$ is a cycle $r$ of degree seven.
By \textup{(D.2.49e1)} write $2r=dR$.  Then
\[
                 2(a'r-ra')=d(a'R-Ra').
 \tag{D.2.49e4}
\]
Thus $2z_L$ is independent of all representatives and top-chain
choices.  Odd orientation changes also leave it fixed because
its order is at most two.

We next compare with the specified actual Moore map, retaining
its entire top cell.  Put $P_5=P^5(2)$ and
$\alpha_P=i_4\nu_4:S^7\to P_5$.  The Hopf--Hilton argument
of D.2.17 gives $4\alpha_P=0$ already in $P_5$, so choose an
extension $\mu_0:M_4\to P_5$.  Let
$c_{2,4}:P^8(2)\to M_4$ be the suspension cofiber map with
bottom degree two and top degree one.  The maps
$\mu_0c_{2,4}$ and $-w=-[i_4,\mathrm{id}_{P_5}]$ have the
same bottom restriction $2\alpha_P$.  Their difference factors
through the top pinch $q_2:P^8(2)\to S^8$ by the Moore
Puppe sequence.  Since $q_4c_{2,4}=q_2$, a top-sphere change
of $\mu_0$ makes
\[
                         \mu_0c_{2,4}=-w
 \tag{D.2.49e5}
\]
as actual homotopy classes.  Here $c_{2,4}$ is a suspension
map, so the comparison respects the coaction without a Hilton
correction.  Composing with $P_5\to Y$ and taking the unique
lift through the six-type quotient defines an actual map
$\mu:M_4\to L$.

The exact natural boundary calculation of D.2.43 gives
\[
 c_{2,4}l_4=\pm\partial(2c_4),\qquad
 \ker(\pi_{14}M_4\to\pi_{14}^{s}M_4)
           =\mathbb Z/8\{\partial c_4\}.
 \tag{D.2.49e6}
\]
The source normalization in this assertion can also be checked
integrally.  The James chain algebra of $\Omega M_4$ is
$T(A_6,B_7)$ with $dB=4A$.  In degree thirteen its only words
are $AB,BA$, and the only degree-fourteen boundary is $d(B^2)$.
Consequently
\[
 H_{13}(\Omega M_4;\mathbb Z)
                    =\mathbb Z/4\{AB-BA\}.
 \tag{D.2.49e7}
\]
There is a stably trivial relative Whitehead generator of
$\pi_{14}M_4$ whose mod-two loop Hurewicz image is the nonzero
commutator.  This is precisely the modulus-four,
bottom-dimension-seven case of
\cite[Theorem~1.7(ii)--(iii), Corollary~1.8, and proof,
pp.~6--9]{EH5ZhuPanComplexes2026}; its order-eight
pinch-kernel presentation is also given in
\cite[Lemma~2.2, p.~6]{EH5ZhuPanClassification2026}.
By \textup{(D.2.49e6)} it is an odd multiple of $\partial c_4$.
Equation \textup{(D.2.49e7)} therefore gives
\[
 h_{\mathbb Z}(\operatorname{adj}\partial c_4)
                  =t(AB-BA),\qquad t\text{ odd}.
 \tag{D.2.49e8}
\]
Only this parity assertion is needed; no numerical choice of
the relative generator is imposed.

The actual loop map $\Omega\mu$ sends these chains to an
allowed pair $a,b$ with $db=4a$.  Naturality,
\textup{(D.2.49e5)}--\textup{(D.2.49e8)} identify the
right side of \textup{(D.2.49e2)} with $\pm2t z_L=2z_L$.
The chain identities \textup{(D.2.49e3)}--\textup{(D.2.49e4)}
make this the same doubled class for every allowed target
choice.  The loop-Hurewicz vanishing just proved now gives
\[
                              2z_L=0.
 \tag{D.2.49e9}
\]
This supplies a chain $W_{14}$ with $dW=2(ab-ba)$ in the
actual loop-chain model.  It does not specify the compatible
higher relation on $W$ or the degree-fifteen Adams--Hilton
differential.  Thus no vanishing of $b_8^2$, of its half-product
pairing on the stable Hurewicz image, or of $h_Ml_4$ in
homotopy is inferred from this conclusion.

\end{proof}
